\documentclass[preprint,12pt]{elsarticle}

\usepackage{amsmath}
\usepackage{amssymb}
\usepackage{booktabs}
\usepackage{enumitem}
\usepackage{graphicx}
\usepackage[hyphens]{url}
\usepackage{hyperref}
\usepackage{makecell}
\usepackage{microtype}
\usepackage{tabularx}
\usepackage{xcolor}
\usepackage{multirow}

\journal{Computers \& Security}

\begin{document}

\begin{frontmatter}

\title{Toward Skill Evolution for\\
Blind Vulnerability Analysis}

\author[inst1]{Authors omitted in working draft}
\address[inst1]{Affiliations omitted in working draft}

\begin{abstract}
Reusable natural-language skills are attractive for blind
vulnerability-analysis agents because they can encode recurring procedures:
identify a dispatcher, narrow to a branch, trace attacker-controlled data,
and localize a sink.  However, a static skill library is not enough.
Across firmware families and bug patterns, a skill that helps one target may
be redundant, ineffective, or even harmful on another.  The core problem is
therefore not only how to write skills, but how to evolve them from
completed runs without polluting the live library seen by later analyses.
We study this problem as \emph{evidence-constrained skill evolution} for
blind vulnerability analysis.  Our setting couples blind remote analysis
with post-run external confirmation: an agent solves a task on an isolated
virtual machine, the local system scores the final answer and checks it
against source predicates, logic harnesses, behavioral oracles, or
sanitizer-backed evidence, and only then may an evolution service propose
and validate candidate skill updates.  We present a split-plane
architecture that separates analysis from confirmation, a run archival
scheme that links selected skills to immutable session segments, a
lightweight run-level feedback interface for candidate generation, and a
publication gate that stages candidate skills before they can enter the
live library.  Empirically, we organize the evidence as a hierarchy rather
than a single benchmark number.  A leakage-controlled 84-run three-condition
rerun shows that earlier near-perfect oracle performance was inflated by
answer-bearing skill content; after cleanup, oracle-skill still outperforms
weak-skill and no-skill while sharply reducing decoy hits.  A larger frozen
153-run necessity study across 17 firmware cases confirms the effect under a
stricter protocol: oracle-skill achieves $46.7\%$ true hits versus
$11.1\%$ for weak-skill and $6.7\%$ for no-skill, while reducing decoy hits
from 28/45 and 24/45 to 2/45.  Complementary studies show that skill
specificity is family-dependent: a distinguishing skill raises Belkin
F9K1122-family accuracy from $20.0\%$ to $66.7\%$, but the same design
principle degrades Tenda FH451-family accuracy from $40.0\%$ to $26.7\%$.
Across 94 archived gate decisions, candidate generation is not the main
bottleneck; defensible publication into the live library is.  Replay-based
gates can reject clearly poor candidates, but they still struggle to
separate genuine improvement from harmless redundancy when the baseline
model already performs well.  The resulting contribution is a
discovery-oriented research scaffold for studying when validated
vulnerability-analysis runs should, and should not, become reusable skills.
\end{abstract}

\begin{keyword}
LLM agent \sep skill evolution \sep vulnerability analysis \sep
external validation \sep blind benchmark \sep evidence-grounded feedback
\end{keyword}

\end{frontmatter}

\section{Introduction}
\label{sec:introduction}

LLM agents increasingly use external memories, reusable instructions, and
natural-language skills to improve over repeated tasks.  Systems such as
Voyager~\cite{voyager} and ExpeL~\cite{expel} show that agents can
accumulate procedural knowledge from interaction and reuse it later.
SkillClaw~\cite{skillclaw} extends this direction to a shared skill
ecosystem, while recent work studies context-to-skill induction
(\cite{ctx2skill}), on-policy skill distillation (\cite{opid}),
environment-grounded skill formation (\cite{spark,skillevolbench}), and
verifier-guided skill refinement (\cite{trace2skill,veriskill}).

This line of work is especially relevant to vulnerability analysis.  Blind
analysis and discovery-oriented auditing repeatedly require the same kinds
of procedural moves: locate a dispatcher, identify a reachable branch,
trace attacker-controlled inputs, follow dataflow toward a sink, and decide
whether the path is exploit-relevant.  Without reusable skills, an agent
must re-learn these procedures on each target, and it often drifts toward a
salient but wrong handler, sink, or vulnerability family.

At the same time, a static skill library is not enough.  Our own firmware
experiments, together with recent work on harmful skills and
self-improvement failures~\cite{skillharmful,misevolution}, suggest that a
skill that helps one case may be ineffective or harmful on another.
Security-specific work has likewise started to evolve reusable playbooks
for auditing~\cite{evohunt} and to construct realistic long-horizon
benchmarks~\cite{secbenchpro}.  The natural next question is therefore not
only why agents need skills, but why those skills must evolve from real
runs rather than remain static templates.

That requirement immediately creates a harder problem.  In a real
vulnerability-analysis workflow, the difficult question is not whether an
LLM or agent can produce a candidate skill revision.  The difficult
question is whether a completed run contains reusable procedural knowledge
that deserves to enter the \emph{live} skill library seen by later runs.
This question is more stringent than ordinary task success.  A run can be
correct yet redundant, correct yet overfit to a single instance, or
correct yet unsafe to generalize.  The engineering surface therefore
includes not only skill retrieval and revision, but also run archival,
evidence collection, candidate staging, and admission control.

Vulnerability analysis is a particularly instructive domain for studying
this problem because it provides something many generic agent tasks do not:
\emph{external, executable confirmation}.  Depending on the case, a claim
can be checked by source predicates, logic harnesses, behavioral oracles,
or sanitizer-backed crashes.  These signals do not prove causal skill
contribution, but they do let us constrain the evolution loop using
evidence that is independent of the model's own narrative.  This suggests a
different research focus from prior work on skill generation or pure
retrieval: rather than asking only how to produce better skills, we can ask
how to govern the transition from a \emph{validated run} to a \emph{live
reusable skill} in a blind vulnerability-analysis workflow.

Our goal is therefore narrower and more operational than a general theory
of skill learning.  We do not claim a new retrieval algorithm, a new
revision model, or a complete solution to causal attribution.  Instead, we
study a full vulnerability-analysis loop in which a blind run executes on a
remote analysis environment, selected skills are injected server-side, the
final answer is scored and externally confirmed on a separate local plane,
the resulting record is transformed into a candidate update, and the update
is published only if it passes an explicit gate.  This framing lets us ask
three concrete questions: how unstable skill effects are in blind security
tasks; whether lightweight post-run feedback is useful enough to drive
candidate generation; and whether publication control, rather than
candidate generation, is the real bottleneck in practice.

This paper makes four contributions:

\begin{enumerate}[leftmargin=*]
    \item \textbf{A skill-evolution framework for blind vulnerability
    analysis.} We formulate the transition from a completed blind run to a
    reusable live skill as a structured lifecycle involving execution,
    confirmation, archival, candidate generation, and publication control.

    \item \textbf{A split-plane architecture for blind analysis and local
    confirmation.} We design a remote-analysis / local-confirmation
    separation that preserves blind execution while still supporting
    post-run evidence processing, session archival, and skill evolution.

    \item \textbf{An evidence-grounded feedback and publication pipeline.}
    We show how externally scored runs can be converted into candidate
    updates while preserving target-identity information, decoy-aware
    diagnostics, and validation-gated publication into the live library.

    \item \textbf{A staged empirical study of why skills help, fail, and
    remain hard to publish.} Through clean reruns, frozen necessity
    studies, family-specific distinguishing experiments, and archived gate
    analysis, we show that skill usage can materially alter vulnerability
    localization quality, but the effect is heterogeneous, family-dependent,
    and not yet sufficient to prove strict necessity in all cases.
\end{enumerate}

The intended scope of this paper is therefore a \emph{framework paper with
initial empirical evidence}.  Our objective is to establish the research
problem, the system structure, and the first round of measurable findings.
We do not claim that autonomous skill improvement has already been solved.
Instead, we argue that the field needs a cleaner way to study when
validated security runs should, and should not, become persistent skills,
and we present a first operational scaffold for that study.

The remainder of the paper is organized as follows.
Section~\ref{sec:problem} formalizes the skill-evolution problem and its
two key subproblems: run-level attribution and safe publication.
Section~\ref{sec:method} presents the framework.
Section~\ref{sec:implementation} describes the prototype.
Section~\ref{sec:evaluation} reports experimental methodology and results.
Section~\ref{sec:discussion} discusses implications.
Section~\ref{sec:related} surveys related work.
Section~\ref{sec:threats} addresses threats to validity.
Section~\ref{sec:conclusion} concludes.

\section{Problem Formulation}
\label{sec:problem}

\subsection{Task Run and Validation Vector}

Let a task run be
\begin{equation}
    r = (x, S_r, \tau, y, E),
\end{equation}
where $x$ is a blind vulnerability-analysis task, $S_r$ is the set of
skills selected and injected for the run, $\tau$ is the agent trajectory,
$y$ is the final answer (predicted files, functions, root cause, and
evidence), and $E$ is externally collected evidence.

We retain a validation vector rather than a single success label:
\begin{equation}
    V(r) = \left(q_{\mathrm{task}},\; q_{\mathrm{confirm}},\;
    a_{r,s}\right).
\end{equation}
Here $q_{\mathrm{task}}$ measures whether the answer localizes and explains
the ground-truth vulnerability; $q_{\mathrm{confirm}}$ represents the
strength of external source or runtime evidence; and $a_{r,s}$ records
\emph{observed relevance} between skill $s$ and the task.

We deliberately call $a_{r,s}$ \emph{observed relevance} rather than
\emph{causal contribution}.  This distinction is the central problem of
this paper.

\subsection{The Skill Attribution Problem}
\label{sec:attribution_problem}

Suppose a run achieves high $q_{\mathrm{task}}$ and high
$q_{\mathrm{confirm}}$: the agent correctly identified the vulnerability
file, function, and root cause, and the result was confirmed by a sanitizer
crash or source predicate.  Can we conclude that the injected skill $s$
causally contributed to this success?

In general, no.  Consider the following decomposition.  Let $y_0$ be the
answer the model would produce \emph{without} skill $s$ (the counterfactual
baseline), and let $y_s$ be the answer \emph{with} skill $s$.  The causal
effect of $s$ on the task outcome is:
\begin{equation}
    \Delta(s) = q_{\mathrm{task}}(y_s) - q_{\mathrm{task}}(y_0).
    \label{eq:causal_effect}
\end{equation}
If $\Delta(s) > 0$, the skill helped.  If $\Delta(s) \leq 0$, it did not
help or was harmful.  However, computing $\Delta(s)$ requires running the
agent twice---once with $s$ and once without---which doubles the
computational cost per attribution.  More importantly, LLM outputs are
stochastic: a single pair of runs may not reflect the true expected
effect, requiring multiple repeated trials for statistical power.

Prior work addresses this through counterfactual ablation.  Differential
comparison between a skill-guided run and a no-skill reference
run~\cite{skillharmful}, and Shapley-style step valuation over skill
traces~\cite{skillshapley,shapley}, both provide stronger attribution
signals than single-run observation.  However, Shapley-style estimation
requires evaluating many subsets of skill steps, making it expensive for
long skills.

The key question we investigate is: \emph{can executable domain evidence
provide a cheaper approximation of $\Delta(s)$ without running
counterfactual ablations?}

\subsection{The Safe Skill Publication Problem}
\label{sec:publication_problem}

When a run is completed and validated, the system may generate a candidate
skill revision $\Delta s$ derived from the run's trajectory and feedback.
The question is whether to publish $\Delta s$ into the live skill library
that subsequent agents will use.

Following the principle that ``evolution optimizes task outcomes rather
than procedure safety''~\cite{misevolution}, we require a publication gate:
\begin{equation}
    \mathrm{Publish}(\Delta s) =
    \mathbb{I}\left[
      \mathrm{EvidenceSufficient}(\Delta s)
      \;\land\;
      \mathrm{FollowupPass}(\Delta s)
    \right].
    \label{eq:gate}
\end{equation}
Here $\mathrm{EvidenceSufficient}$ checks whether the originating run had
adequate evidence to justify a skill update, and $\mathrm{FollowupPass}$
runs a follow-up validation with the candidate skill to check whether it
maintains or improves performance.

The structural challenge is in $\mathrm{FollowupPass}$.  In a
\emph{replay-only} paradigm, the follow-up runs the candidate skill on the
same or similar task and checks whether the result passes validation.  But
this only measures $q_{\mathrm{task}}(y_{\Delta s})$, not
$\Delta(\Delta s) = q_{\mathrm{task}}(y_{\Delta s}) -
q_{\mathrm{task}}(y_0)$.  Without the counterfactual baseline $y_0$, the
gate cannot distinguish ``the skill helped'' from ``the model would have
succeeded anyway.''  When the gate requires \emph{strict improvement}
(mean score with candidate $>$ mean score without candidate), this leads to
systematic rejection of candidates that are not harmful but merely
redundant---a high false-negative rate.

\subsection{Research Hypotheses}

The study is guided by three hypotheses:

\begin{description}
    \item[H1: Clean skill effect.] Under leakage-controlled blind
    execution, stronger domain skills should increase true-hit rate and
    suppress decoy-hit rate relative to weak-skill and no-skill controls,
    but the effect remains case-dependent.
    \item[H2: Specificity boundary.] More specific family-level guidance can
    recover part of the residual error left by generic skills, but only for
    families that expose stable discriminative cues.
    \item[H3: Gate limitation.] A replay-only publication gate suffers a
    structural false-negative limitation that cannot be resolved by
    parameter tuning alone.
\end{description}

\section{Evidence-Constrained Skill Evolution Framework}
\label{sec:method}

We propose a framework that connects blind task execution, skill injection
auditing, external evidence confirmation, evidence-conditioned feedback,
and gated candidate publication.  The framework is designed around a simple
principle: in vulnerability analysis, post-run evidence should constrain
both what we infer about a skill and what we allow the system to retain for
future use.  The goal is not to make the evolution service omniscient.  The
goal is to ensure that the path from a completed run to a persistent live
skill is explicit, inspectable, and conservative enough for security work.

\subsection{Leakage-Controlled Blind Execution}
\label{sec:blind}

Each evaluation begins from a benchmark specification containing private
ground truth and a separately rendered blind prompt.  The agent receives
the target workspace and the analysis objective, but not the answer fields,
confirmation scripts, or solution-bearing artifact names.  The benchmark
controller retains the expected file, function, root cause, and
confirmation policy.

The analysis runs on a remote virtual machine, separated from the
benchmark ground truth and from the local service that evaluates the
result.  This separation is methodologically important: it prevents
inadvertent answer leakage through shared filesystems or local environment
artifacts, and it mirrors real-world deployment where the analyst's
environment is isolated from the vulnerability database.

\subsection{Auditable Skill Injection}
\label{sec:injection}

For each run, the system records the selected skill names, the injected
prompt content hash, the request session identifier, and the immutable
session segment consumed by the evolution process.

This recording addresses a prerequisite for studying skill evolution.
Any claim about a skill's effect is uninterpretable if the intended skill
cannot be linked to a specific request segment, or if the injected content
cannot be reconstructed after the run.  Injection auditing therefore does
not solve attribution, but it narrows the question from ``was the skill
present?'' to the more meaningful ``what procedural prior was exposed to
the model in this specific run, and how should that exposure be archived
for later judgment?''

\subsection{Multi-Evidence External Confirmation}
\label{sec:confirmation}

The agent's output is evaluated along two complementary dimensions.  First,
\emph{structured scoring} checks the predicted vulnerability identity,
file, function, root cause, and cited evidence against ground truth.
Second, \emph{case validators} test claims against the target environment.

We classify evidence into four categories, each making a different claim
about the vulnerability:

\begin{itemize}[leftmargin=*]
    \item \textbf{Source-backed evidence:} verifies a concrete code
    predicate or boundary relation (e.g., an unsafe function call exists
    at the predicted location).
    \item \textbf{Behavior-backed evidence:} reproduces an expected program
    path or output condition (e.g., the binary processes the input and
    reaches the predicted function).
    \item \textbf{Logic-backed evidence:} exercises a controlled harness
    for the vulnerable state transition (e.g., a test input triggers the
    expected code path without a full crash).
    \item \textbf{Sanitizer-backed evidence:} reproduces a memory-safety
    failure with relevant runtime diagnostics (e.g., AddressSanitizer
    reports a buffer overflow at the predicted location).
\end{itemize}

The distinction matters because different evidence classes support
different claims about the analysis.  A sanitizer crash is stronger
evidence for a memory-safety vulnerability than a source-code predicate
alone.  The framework records the accepted evidence class rather than
mapping every passing script to the same notion of ``confirmed.''

\subsection{Evidence-Conditioned Feedback}
\label{sec:feedback}

After validation, the system constructs feedback at the run--skill level.
This feedback is an operational interface for evolution rather than a
standalone causal estimator.  Its role is to summarize the relation between
task outcome, external confirmation, and observed skill relevance in a form
that can drive candidate generation and later gate decisions.  In that
sense, it is a low-cost approximation to the latent quantity $\Delta(s)$ in
Equation~\ref{eq:causal_effect}, but we do not present it as a substitute
for full counterfactual attribution.

A skill receives one of four statuses:

\begin{description}[style=nextline,leftmargin=2.4cm]
    \item[Supported ($\hat{\Delta} > 0$):] the task is externally validated
    and the skill is aligned with the analysis procedure (i.e., the skill's
    described steps match the agent's actual trajectory).  This is
    positive evidence for the skill, but not proof of causality.
    \item[Mismatched ($\hat{\Delta} \approx 0$):] the task may be correct,
    but the skill is unrelated to the target or method.  The model likely
    succeeded on its own; the skill was redundant.
    \item[Contradicted ($\hat{\Delta} < 0$):] the skill appears to drive a
    false localization, unsupported claim, or failed confirmation.  The
    skill was actively harmful.
    \item[Unknown:] the available trace cannot distinguish contribution
    from coincidence.
\end{description}

The key idea is practical rather than metaphysical.  Vulnerability-analysis
workflows already produce task scores, evidence classes, and auditable
records of the selected skill and final answer.  We reuse these signals to
derive an attribution-oriented label at the cost of one run, whereas
counterfactual ablation requires at least two runs and Shapley-style
valuation~\cite{skillshapley} can be exponentially expensive in the number
of skill steps.  This makes the label useful as a triage signal for
evolution even when it remains imperfect as scientific evidence of causal
contribution.

However, this approximation has lower precision.  A ``supported'' skill may
be supported only by coincidence: the model might have succeeded regardless.
A ``mismatched'' skill might have subtly helped in ways the alignment check
cannot detect.  We empirically evaluate this precision trade-off in
Section~\ref{sec:results}.

\subsection{Candidate-Level Publication Gate}
\label{sec:gate}

Evolution produces a candidate revision $\Delta s$ rather than directly
overwriting the live skill.  Publication is gated by
Equation~\ref{eq:gate}.  In our prototype, follow-up validation replays the
candidate on the same or a closely related task and checks whether the
candidate remains consistent with the stored evidence.

The most conservative form of this gate is a \emph{strict improvement}
policy: the candidate must achieve a mean validation score that exceeds a
baseline score by a margin.  This design is motivated by the concern that
successful-but-unsafe or merely coincidental experiences should not be
persisted as reusable skills.

However, as formalized in Section~\ref{sec:publication_problem}, a
replay-only gate has a structural limitation.  The follow-up measures
$q_{\mathrm{task}}(y_{\Delta s})$ but not
$\Delta(\Delta s) = q_{\mathrm{task}}(y_{\Delta s}) -
q_{\mathrm{task}}(y_0)$.  When the baseline is high (the model can already
solve the task), strict improvement is nearly impossible to establish
because the candidate's replay score is bounded by the same ceiling.  We
provide empirical evidence of this limitation in
Section~\ref{sec:results_gate}.

\subsection{Design Rationale: Live, Candidate, and Share Separation}

Live skills, candidate revisions, and shared runtime state are stored in
separate directories.  This separation is not merely a
software-organization choice; it is part of the experimental design.  It
ensures that candidate skills can be generated, inspected, replayed, and
rejected without silently changing the library seen by the next run.

\section{Prototype Realization}
\label{sec:implementation}

We realize the framework as a research extension to
SkillClaw~\cite{skillclaw}.  The original system already provides request
proxying, skill retrieval and injection, session capture, skill sharing,
and an evolution service.  Our extension does not replace these core
services with a new retrieval or revision algorithm.  Instead, it adds the
elements needed to study the run-to-library transition in a security
setting: leakage-controlled benchmark cases, remote blind execution,
heterogeneous vulnerability validators, candidate staging, and a structured
handoff from finalized runs to validation-gated evolution.

\subsection{Analysis Plane and Evidence Plane}

The runtime is divided into two planes.  The \emph{analysis plane} consists
of an LLM agent client on a remote VM that analyzes the target through the
SkillClaw proxy.  The \emph{evidence plane} is a local service that retains
private ground truth, finalizes the run record, executes validators, and
sends a closed session/feedback pair to the evolution service.  This split
preserves the blind-evaluation setting while still allowing rich local
post-processing and confirmation.

Client-provided session identifiers are preserved when available, and long
sessions are archived as incrementally consumable immutable segments.  This
ensures that later interaction can extend a session without rewriting
feedback already consumed by the evolver.

\subsection{Benchmark Cases}

We define benchmark cases in JSON specifications containing: the blind
workspace path, the analysis objective (rendered without CVE identifiers or
solution-bearing names), private ground truth (expected files, functions,
root cause, required evidence terms), and confirmation policies.

The repository-level benchmark library currently contains 27 case
specifications, but they are not all used in the same way.  We organize
them into three evidence tiers:
\begin{itemize}[leftmargin=*]
    \item \textbf{Primary controlled firmware cases:} blind embedded HTTPd
    command-injection, overflow, and path-traversal tasks analyzed through
    binary-level reasoning on a remote VM.  The main quantitative results
    in this paper use a frozen 17-case subset from this tier.
    \item \textbf{Open-source confirmation cases:} parser vulnerabilities
    in giflib, libxml2, tcpdump, libarchive, and Exiv2, where source code,
    sanitizer runs, and logic-backed validators provide stronger external
    confirmation.  These cases are valuable for validator and gate
    development, but they are not folded into the frozen 153-run firmware
    denominator.
    \item \textbf{Supporting diagnostic cases:} earlier development-time
    or mechanism-isolation cases, including focused CGI profile ablations
    and decoy-heavy family studies, retained to explain scoring drift,
    skill structure, and failure modes.
\end{itemize}

This separation matters methodologically.  Firmware cases emphasize blind
target localization under binary-only reasoning, whereas the open-source
cases emphasize richer confirmation and validator design.  Mixing them into
a single denominator would blur the distinction between skill effect and
evidence availability.

\subsection{Scoring}

Each case defines a scoring rubric with four check dimensions:
\begin{itemize}[leftmargin=*]
    \item \textbf{File localization:} does the agent identify the correct
    vulnerable file?
    \item \textbf{Function localization:} does the agent identify the
    correct vulnerable function?
    \item \textbf{Root cause:} does the agent describe the root cause with
    sufficient technical detail?
    \item \textbf{Evidence:} does the agent cite concrete evidence
    (function names, string references, code patterns)?
\end{itemize}
Each dimension is scored against private case metadata and aggregated onto a
0--10 scale through a weighted rubric.  The rubric is deliberately coarse:
it is intended to compare blind runs at the task level, not to certify
fine-grained semantic equivalence between two explanations.

\subsection{Evolution Service}

The evolution service consumes the feedback bundle and generates candidate
skill revisions using an LLM-based workflow engine.  Candidates are
content-deduplicated by SHA-256 hash before entering the gate queue.  The
gate runs follow-up validation and enforces the publication policy defined
in Section~\ref{sec:gate}.

\subsection{LLM Configuration}

All controlled experiments reported as primary quantitative evidence use the
same backend, \texttt{deepseek-v4-flash}, accessed through the local
CC-Switch proxy.  This does not remove output stochasticity, but it avoids
cross-model capability differences within the main comparisons.  Earlier
development-time diagnostic runs span older runtime snapshots and are used
only as supporting evidence for protocol repair, not as the main basis for
the paper's claims.

\section{Experimental Methodology}
\label{sec:evaluation}

\subsection{Research Questions}

\begin{description}
    \item[RQ1.] Under leakage-controlled blind execution, does skill
    strength measurably change true-hit rate and decoy-hit rate?
    \item[RQ2.] When generic skills still fail, can family-specific
    distinguishing skills recover the target function, and under what
    boundary conditions?
    \item[RQ3.] After the loop is operational, what does the publication
    gate actually guarantee, and what ambiguity remains between genuine
    improvement and harmless redundancy?
\end{description}

\subsection{Benchmark Cases}

The primary benchmark is a frozen set of 17 embedded firmware cases drawn
from HTTP handler and CGI vulnerabilities.  The controlled studies are run
through blind binary-level reasoning on a remote VM, with neither CVE
identifiers nor solution-bearing artifact names exposed to the agent.  The
cases span command-injection, buffer-overflow, and path-traversal
localization tasks across several firmware families.

\begin{table}[t]
\centering
\caption{Families represented in the controlled benchmark.}
\label{tab:cases}
\footnotesize
\begin{tabular}{l c l}
\toprule
Family / device line & Cases & Main challenge \\
\midrule
Tenda F1202 & 2 & command injection and credential-related overflow \\
Tenda F453 & 4 & multiple HTTP handlers with strong \texttt{formexeCommand} decoys \\
Tenda F456 & 1 & command injection near the same Tenda sink cluster \\
Belkin F9K1122 & 4 + 2 near-pair cases & family-internal handler confusion around \texttt{formWISP5G} and settings handlers \\
Tenda FH451 & 5 & dense handler clusters with overlapping token patterns \\
i12 & 1 & path traversal versus adjacent access-control handlers \\
\bottomrule
\end{tabular}
\end{table}

Two additional earlier studies are retained as supporting diagnostics rather
than merged into the main benchmark: a 16-run firmware CGI profile ablation
over \texttt{firmware2-wireless} and \texttt{firmware2-login}, and a
16-run Tenda F453 blind diagnostic study that exposed score/feedback
misalignment.  These runs remain useful for mechanism analysis, but they
were produced under earlier development-time protocols and are therefore
not treated as the primary evidence base.

\subsection{Conditions and Controls}

The main evaluation uses three skill conditions:
\begin{itemize}[leftmargin=*]
    \item \textbf{No-skill:} no skill is injected.
    \item \textbf{Weak-skill:} a generic, non-answer-bearing procedural
    skill is injected (e.g., broad CGI or ELF triage guidance).
    \item \textbf{Oracle-skill:} a case-family-specific skill is injected,
    still audited to remove direct answer leakage.
\end{itemize}

We use two controlled protocols:
\begin{itemize}[leftmargin=*]
    \item \textbf{Clean rerun protocol:} 14 cases $\times$ 3 conditions
    $\times$ 2 rounds = 84 runs.  This protocol was introduced to verify
    that earlier near-perfect oracle performance was caused by
    answer-bearing skill content.
    \item \textbf{Frozen necessity protocol:} 17 cases $\times$ 3
    conditions $\times$ 3 rounds = 153 runs.  This is the main quantitative
    study reported in the paper.
\end{itemize}

For the frozen protocol, two F9K1122 near-pair handlers
(\texttt{formSetSystemSettings} and \texttt{formWlanSetup}) are reported
separately rather than folded into the per-function denominator, because
generic oracle wording cannot reliably distinguish them under token-level
blind reasoning.  This prevents the aggregate statistics from overstating
or understating skill effect on separable targets.

\subsection{Complementary Targeted Profile Ablations}

We retain four complementary studies for mechanism diagnosis:
\begin{itemize}[leftmargin=*]
    \item a 16-run firmware CGI profile ablation comparing \emph{none},
    \emph{wrong}, \emph{relevant}, and \emph{seed} skills on
    \texttt{firmware2-wireless} and \texttt{firmware2-login};
    \item a 16-run F453 blind diagnostic study that revealed repeated drift
    from \texttt{formWriteFacMac} to the decoy handler
    \texttt{formexeCommand};
    \item a 15-run Belkin F9K1122 family-specific distinguishing study;
    \item a 15-run Tenda FH451 family-specific distinguishing study.
\end{itemize}

These studies are not merged into the frozen aggregate.  They answer more
focused questions about skill structure, decoy attraction, and family-level
separability.

\subsection{Measurements}

We report:
\begin{itemize}[leftmargin=*]
    \item \textbf{Task score:} a 0--10 score decomposed into file,
    function, root-cause, and evidence sub-scores.
    \item \textbf{Outcome class:} \textbf{YES} (true target hit),
    \textbf{DECOY} (landed on a dangerous but wrong target), or
    \textbf{NO} (miss).
    \item \textbf{Feedback metadata:} including target-aware flags such as
    \texttt{function\_identity\_miss} when the run is close but lands on the
    wrong handler.
    \item \textbf{Gate outcome:} published, rejected, or pending, together
    with the gate mode and baseline relation when available.
\end{itemize}

\subsection{Evidence Hierarchy and Reporting Policy}

Not all experimental artifacts produced during system development are
equally suitable as paper evidence.  We therefore report results under an
explicit hierarchy:
\begin{itemize}[leftmargin=*]
    \item \textbf{Primary quantitative evidence:} the 84-run clean rerun,
    the 153-run frozen necessity study, and the targeted F9K1122/FH451
    family-specific distinguishing experiments.
    \item \textbf{Supporting diagnostic evidence:} earlier firmware2 CGI
    profile ablations and the F453 blind diagnostic study, which are used
    to interpret failure modes, decoy attraction, and scoring/gate
    behavior.
    \item \textbf{Retired or non-primary evidence:} development snapshots
    later found to include answer-bearing skill content, contaminated
    prompts, or benchmark definitions that were not yet protocol-clean.
\end{itemize}

This policy is important for two reasons.  First, it prevents earlier
development-time wins from being mistaken for stable benchmark results.
Second, it keeps negative findings visible: studies that exposed leakage,
mis-scoring, or invalid case composition are not deleted, but they are no
longer allowed to dominate the paper's main quantitative claims.

\subsection{Statistical Analysis}

Our analysis is primarily descriptive.  The clean rerun protocol provides
$N=2$ rounds per case-condition and the frozen protocol provides $N=3$.
This is enough to expose large directional effects and failure modes, but
not enough for strong per-case significance claims.  We therefore emphasize
aggregate rates, case-level breakdowns, and protocol transparency rather
than inferential statistics.

\section{Results}
\label{sec:results}

\subsection{RQ1: Clean Controlled Studies Show Real but Incomplete Skill Effects}
\label{sec:results_rq1}

Table~\ref{tab:clean_frozen_results} summarizes the two main controlled
protocols.

\begin{table}[t]
\centering
\caption{Leakage-controlled three-condition results.}
\label{tab:clean_frozen_results}
\footnotesize
\begin{tabular}{l l c c c c c}
\toprule
Protocol & Condition & $n$ & YES & DECOY & NO & Mean score \\
\midrule
Clean rerun (84 runs) & no-skill & 28 & 1 (3.6\%) & 9 & 18 & 4.38 \\
Clean rerun (84 runs) & weak-skill & 28 & 4 (14.3\%) & 12 & 12 & 4.34 \\
Clean rerun (84 runs) & oracle-skill & 28 & 10 (35.7\%) & 1 & 17 & 6.32 \\
\midrule
Frozen necessity (153 runs) & no-skill & 45 & 3 (6.7\%) & 28 & 14 & 3.92 \\
Frozen necessity (153 runs) & weak-skill & 45 & 5 (11.1\%) & 24 & 16 & 3.99 \\
Frozen necessity (153 runs) & oracle-skill & 45 & 21 (46.7\%) & 2 & 22 & 6.20 \\
\midrule
\multicolumn{7}{p{0.94\linewidth}}{\footnotesize Frozen-necessity denominator excludes two F9K1122 near-pair handlers that are reported separately under the protocol.} \\
\bottomrule
\end{tabular}
\end{table}

\textbf{Finding 1: protocol cleanup was necessary.}
The clean-rerun study was introduced because earlier oracle runs had
produced near-perfect performance under development-time prompts.  After
removing answer-bearing skill content, oracle accuracy fell to $35.7\%$.
This confirms that the older oracle numbers were not acceptable as
scientific evidence.  The current paper therefore treats the clean and
frozen studies, rather than earlier development snapshots, as the main
quantitative basis.

\textbf{Finding 2: skill changes target selection, not just total score.}
Under the frozen protocol, oracle-skill reaches $46.7\%$ true hits, versus
$11.1\%$ for weak-skill and $6.7\%$ for no-skill.  More importantly,
oracle-skill reduces decoy hits from 28 and 24 down to 2.  In other words,
the strongest observable effect is not only a score increase, but a sharp
reduction in being pulled toward the wrong dangerous handler.

\textbf{Finding 3: the effect is real but not yet strict necessity.}
Oracle-skill still leaves 22/45 runs in the frozen denominator as NO.
Several cases remain resistant even under oracle guidance, including
\texttt{f1202-credential-overflow}, \texttt{f453-qossetting-overflow}, and
\texttt{fh451-WrlclientSet-overflow}.  The correct claim at this stage is
therefore not ``the skill is always necessary,'' but rather ``skill
strength materially shifts the probability of landing on the true
target.''

\textbf{Finding 4: some families are already close to solvable under the
current oracle formulation.}
For example, the frozen study shows
\texttt{f1202-cmdinject} moving from $1/3$ true hits under no-skill and
weak-skill to $3/3$ under oracle-skill; \texttt{f453-overflow} from
$0/3,0/3$ to $3/3$; and \texttt{f456-cmdinject} from $1/3,1/3$ to $3/3$.
These are not universal wins, but they demonstrate that the current skill
format can already create strong target-selection differences on some
families.

\subsection{RQ2: Skill Specificity Helps Only When the Family Is Separable}
\label{sec:results_rq2}

We next ask why oracle still fails on many frozen cases.  The answer from
the supporting studies is that ``more specific'' is not universally better;
it depends on whether the firmware family exposes stable intra-family
signals.

\begin{table}[t]
\centering
\caption{Supporting diagnostic studies on skill structure and family-level specificity.}
\label{tab:diagnostic_results}
\footnotesize
\begin{tabular}{l l l}
\toprule
Study & Conditions & Result \\
\midrule
Firmware2-wireless & none / wrong / relevant / seed & 2.67 / 3.00 / 3.00 / \textbf{4.00} \\
Firmware2-login & none / wrong / relevant / seed & 4.33 / 4.50 / \textbf{5.00} / \textbf{5.00} \\
Belkin F9K1122 & generic oracle / distinguishing oracle & 20.0\% $\rightarrow$ \textbf{66.7\%} \\
Tenda FH451 & generic oracle / distinguishing oracle & 40.0\% $\rightarrow$ 26.7\% \\
\bottomrule
\end{tabular}
\end{table}

\textbf{Firmware2 profile ablations: structure matters.}
The earlier firmware2 CGI ablation does not directly measure necessity, but
it shows that a branch-first seed skill can outperform a broader
descriptive live skill on \texttt{firmware2-wireless} while tying it on
\texttt{firmware2-login}.  Topical relevance alone is therefore not enough;
the internal organization of the skill can change how the model traverses a
CGI dispatcher.

\textbf{F9K1122: specificity can rescue a family.}
In the frozen study, generic oracle still leaves the family vulnerable to a
\texttt{formWISP5G} attractor.  The distinguishing F9K1122 skill adds
family-level cues such as adjacent symbol and parameter-pattern
distinctions, raising the correct rate from $20.0\%$ to $66.7\%$.  This is
the clearest positive example in the current repository that a more
specific skill can materially recover target localization.

\textbf{FH451: specificity can also fail.}
Applying the same design idea to FH451 lowers the family-level result from
$40.0\%$ to $26.7\%$.  The lesson is not that distinguishing skills are
wrong, but that this style of specificity only helps when the family
actually exposes usable discriminative cues.  FH451 appears to have denser
handler overlap and weaker token separation than F9K1122.

\textbf{F453: wrong-target high scores motivated target-aware feedback.}
The earlier 16-run F453 diagnostic study produced only 2 exact hits among
15 valid runs, while 10 runs drifted to the wrong but dangerous handler
\texttt{formexeCommand}.  Many of those wrong-path runs still scored 8.0
because they hit the correct binary and sink family.  This is exactly the
engineering failure mode that later motivated target-aware metadata such as
\texttt{function\_identity\_miss} in the feedback and gate pipeline.

\subsection{RQ3: Replay-Only Gate Has a Structural Limitation}
\label{sec:results_gate}

We analyze 94 gate decisions collected across all experiments
(Table~\ref{tab:gate_stats}).

\begin{table}[t]
\centering
\caption{Gate decision summary across all experiments.}
\label{tab:gate_stats}
\small
\begin{tabular}{lr}
\toprule
Metric & Count \\
\midrule
Total gate decisions & 94 \\
Published & 12 \\
Rejected & 82 \\
Published under earlier score-only gate logging & 10 \\
Published under baseline-aware non-inferiority logging & 2 \\
\bottomrule
\end{tabular}
\end{table}

The archived decision history spans two gate regimes.  Earlier decisions
record only the candidate-side validation score and final decision.  Later
decisions additionally record a baseline mean, gate mode, and effective
threshold.  Among the latter, the two published examples were accepted
under a non-inferiority criterion, with $(0.7, 0.4)$ and $(0.6, 0.6)$ for
(candidate mean, baseline mean), respectively.

\textbf{Operationally, the loop is now closed.}
The repository already contains real candidate $\rightarrow$ rerun
$\rightarrow$ gate $\rightarrow$ publish examples, and the gate-settle
audit further repaired 21 stuck pending jobs down to 0, set the default
\texttt{publish\_mode} to \texttt{validated}, and added unit coverage for
the settle logic.  This matters because the current question is no longer
whether the system can move a candidate through the pipeline; it can.  The
question is what that publication actually proves.

\textbf{Root cause analysis.}
The key issue is not simply that many candidates are poor.  Rather, the
gate often lacks a defensible basis for deciding whether a candidate skill
was \emph{necessary}.  In a replay-only setting, a candidate can reproduce
the same validated answer as the baseline while providing no clear evidence
that it changed the outcome.  Under a strict-improvement interpretation,
such a candidate is rejected as redundant; under a non-inferiority
interpretation, the same candidate may be retained as harmless but
potentially useful.  The decision therefore depends as much on the gate's
philosophy as on the candidate's intrinsic quality.

This is the structural limitation formalized in
Section~\ref{sec:publication_problem}.  The gate's
$\mathrm{FollowupPass}$ function computes $q_{\mathrm{task}}(y_{\Delta s})$
but only imperfectly captures $q_{\mathrm{task}}(y_0)$.  When the model can
already solve the task without the skill, a replay-only gate struggles to
separate three cases: genuine improvement, harmless redundancy, and
spurious success.  This ambiguity is exactly why publication remains the
hardest part of the loop.

\textbf{Interpretation.}
Recent work on skill misevolution~\cite{misevolution} proposes SafeEvolve, a wrapper that repairs unsafe
content and governs subsequent reuse.  Our study highlights a closely
related but distinct problem: even when content appears safe, the system
may still be unable to decide whether the skill deserves to persist.  In
other words, safety and usefulness are separable questions, and a realistic
evolution loop must reason about both.

\subsection{Summary of Findings}

\begin{table}[t]
\centering
\caption{Summary of findings mapped to hypotheses.}
\label{tab:findings}
\small
\begin{tabularx}{\textwidth}{lXl}
\toprule
Hypothesis & Finding & Status \\
\midrule
H1: Clean skill effect & Under clean and frozen protocols, oracle-skill sharply outperforms weak/no skill and suppresses decoys, but many cases still remain unsolved & Supported \\
H2: Specificity boundary & F9K1122 benefits strongly from distinguishing skills, whereas FH451 does not & Supported \\
H3: Gate limitation & The loop can publish candidates, but replay-only evidence still cannot cleanly separate improvement from redundancy & Supported \\
\bottomrule
\end{tabularx}
\end{table}

\section{Discussion}
\label{sec:discussion}

\subsection{The Evidence--Attribution Gap}

Our results reveal what we call the \emph{evidence--attribution gap}:
executable domain evidence can verify \emph{what} the agent found (the
vulnerability is real, the localization is correct), but it cannot
directly verify \emph{why} the agent found it (was the skill causally
responsible?).

This gap has two dimensions:
\begin{enumerate}[leftmargin=*]
    \item \textbf{Positive evidence is not causal evidence.} A run that
    passes source-backed validation and has a ``supported'' feedback status
    may have succeeded because of the skill, because of the model's prior
    knowledge, or because of both.  The four-level feedback reduces this
    ambiguity by checking skill--trajectory alignment, but it cannot
    eliminate it.
    \item \textbf{Absence of evidence is not evidence of absence.} A
    ``mismatched'' skill may have subtly influenced the agent's reasoning
    in ways that the alignment check cannot detect.  For example, a skill
    might prime the agent to look at certain binary sections, indirectly
    leading to the correct answer even if the skill's specific steps were
    not followed.
\end{enumerate}

Bridging this gap requires either counterfactual ablation (expensive but
precise) or a probabilistic attribution model that estimates $P(\text{success}
\mid \text{skill}) - P(\text{success} \mid \text{no skill})$ from observed
correlations (cheaper but requires more data).  We leave the latter as
future work.

\subsection{Implications for Skill Library Design}

Our finding that the same skill can be helpful, neutral, or harmful
depending on the case has a direct implication: \emph{skill libraries need
case-sensitive selection and case-sensitive evaluation.}  A globally
plausible skill may still degrade a specific analysis if its procedural
biases do not fit the target.  This is consistent with
prior observations that skill-induced failures are strongly
case-specific~\cite{skillharmful}.

This suggests that retrieval mechanisms should consider not only textual
skill--task similarity, but also uncertainty about the model's unaided
capability.  If the baseline is already strong, injection may add noise,
encourage over-commitment to the wrong pattern, or simply consume prompt
budget without benefit.

\subsection{Toward a Counterfactual Gate}
\label{sec:counterfactual_gate}

The structural limitation of the replay-only gate
(Section~\ref{sec:results_gate}) suggests two directions for improvement:

\textbf{Scheme A: Counterfactual ablation gate.} For each candidate, run
the task both with and without the candidate skill, and require strict
improvement ($\bar{y}_{WS} > \bar{y}_{NS}$).  This directly measures
$\Delta(\Delta s)$ but doubles the validation cost and requires
statistical power (multiple rounds).

\textbf{Scheme B: Non-inferiority gate.} Instead of requiring strict
improvement, require that the candidate is \emph{not worse} than the
baseline by more than a margin $\epsilon$: $\bar{y}_{WS} \geq
\bar{y}_{NS} - \epsilon$.  This relaxes the gate to allow redundant (but
not harmful) candidates, at the cost of potentially publishing skills that
do not help.

A combined approach could use Scheme B for initial publication (allowing
the skill to enter the live library) and Scheme A for retention decisions
(removing skills that do not demonstrate improvement over a longer
horizon).  More broadly, this suggests that publication and retention may
need to be decoupled: the system can admit a candidate as provisionally
safe while still demanding stronger evidence before treating it as
established expertise.

\subsection{Comparison with Concurrent Work}

Table~\ref{tab:comparison} compares our approach with concurrent work on
agent skills.

\begin{table}[t]
\centering
\caption{Comparison with concurrent work on agent skills.}
\label{tab:comparison}
\footnotesize
\begin{tabular}{l c c c c c}
\toprule
& \multicolumn{2}{c}{Attribution} & \multicolumn{2}{c}{Publication} & \\
\cmidrule(lr){2-3} \cmidrule(lr){4-5}
Approach & Method & Cost & Method & Safety & Domain \\
\midrule
Ours & 4-level evidence & 1 run & Replay gate & Conservative & Vuln. analysis \\
SkillShapley~\cite{skillshapley} & Shapley value & $O(2^n)$ & N/A & N/A & General \\
Skill Harmful~\cite{skillharmful} & Differential & 2 runs & N/A & N/A & General \\
Misevolution~\cite{misevolution} & N/A & N/A & SafeEvolve & Repair-based & General \\
\bottomrule
\end{tabular}
\end{table}

Our approach is distinctive in its use of executable domain evidence as a
post-run attribution signal.  This gives it a cost advantage over explicit
ablation, but also a narrower claim: we obtain an auditable approximation,
not a full causal decomposition.

\subsection{Limitations of Static Analysis}

The firmware cases in our A/B study are analyzed through binary-level
static reasoning without a live runtime environment.  Consequently, the
available evidence is weaker than sanitizer-backed confirmation in the
open-source cases.  This likely lowers the precision of the feedback signal
and should be interpreted as a limitation of the present study, not as a
property of the framework in general.

\section{Related Work}
\label{sec:related}

\subsection{Skill Induction, Libraries, and Evolution}

Voyager~\cite{voyager}, ExpeL~\cite{expel}, Reflexion~\cite{reflexion},
and MemGPT~\cite{memgpt} establish the broader agent-learning context in
which trajectories, memories, and reusable instructions are accumulated
across tasks.  SkillClaw~\cite{skillclaw} makes this logic explicit at the
level of a shared natural-language skill library, while
Ctx2Skill~\cite{ctx2skill} and OPID~\cite{opid} study, respectively,
context-to-skill induction and trajectory-derived skill distillation.
These works show that skill formation and reuse are already central to LLM
agent research.  Our paper does not compete by proposing yet another
generic skill generator.  Instead, it focuses on a narrower systems
question: in a vulnerability-analysis deployment, how should a completed
run be transformed---or refused transformation---into a live reusable
skill?

\subsection{Environment- and Verifier-Grounded Skill Refinement}

Recent work increasingly grounds skill refinement in environmental or
verifier feedback rather than in free-form self-reflection alone.
SPARK~\cite{spark} emphasizes \emph{evidence over plans}, preferring
environment-grounded revisions to purely narrative improvement.
Trace2Skill~\cite{trace2skill} distills trajectory-local lessons into
longer-term skills using verifier feedback, and
VeriSkill~\cite{veriskill} studies self-evolution specifically for program
verification skills.  SkillEvolBench~\cite{skillevolbench} further turns
episodic-to-procedural evolution into an explicit benchmark problem.

Our work is aligned with this trend, but differs in scope and target.  We
do not study general verifier-backed skill improvement in the abstract.
We study the security-specific transition from \emph{blind remote
vulnerability-analysis runs} to \emph{candidate skills that may alter a
live library}.  In this setting, the central issue is not only whether
verifier signals exist, but how they participate in archival, candidate
staging, and publication control.

\subsection{Skill Failure, Attribution, and Governance}

Skill-induced failures~\cite{skillharmful} demonstrate that apparently
relevant skills can actively harm agent performance by steering execution
onto the wrong procedure.  Skill misevolution~\cite{misevolution} shows
that self-improving loops can preserve unsafe behavior as persistent
skills.  Adversarial skill publishing~\cite{detourhijack} further exposes
how untrusted skills can manipulate downstream behavior.

These works motivate our emphasis on post-run governance.  In our setting,
retrieval and generation are only the front half of the problem.  The back
half is deciding whether a candidate skill deserves admission into the live
library.  Our gate analysis complements prior work by showing that even
when candidate generation is available, publication remains difficult
because replay-only evidence is often insufficient to separate improvement
from redundancy.

\subsection{Skill Attribution and Valuation}

SkillShapley~\cite{skillshapley} formulates skill-step attribution through
Shapley-value estimation~\cite{shapley}, yielding a principled but costly
counterfactual approach.  Differential no-skill / with-skill
comparison~\cite{skillharmful} provides a cheaper alternative at the run
level.  Our feedback interface belongs to this family of post-run
assessment methods, but with a deliberately narrower claim: it is a
triage-oriented operational signal, not a full causal estimator.  Its
purpose is to support candidate generation and publication decisions in a
security workflow where full ablations are expensive and often delayed.

\subsection{Security Agents, Playbooks, and Benchmarks}

On the security side, PentestGPT~\cite{pentestgpt} and subsequent studies
on LLM-based vulnerability analysis~\cite{gpt4vuln,aslr} focus on what
models can discover or explain.  More recent work such as
EvoHunt~\cite{evohunt} begins to evolve reusable playbooks for agentic
security auditing, while SEC-bench Pro~\cite{secbenchpro} raises the bar on
realistic long-horizon software-security evaluation.  Our work is adjacent
to both directions.  Like EvoHunt, we care about reusable procedural
knowledge in security tasks; like SEC-bench Pro, we care about realistic
task construction and leakage control.  The difference is that we center
the run-to-library transition itself: blind execution, external
confirmation, candidate staging, and live-skill admission are the primary
objects of study.

\section{Threats to Validity}
\label{sec:threats}

\paragraph{Sample size.}
The main controlled studies are larger than the early six-case A/B stage,
but per-case repetition is still limited: the clean rerun uses $N=2$ rounds
per condition and the frozen protocol uses $N=3$.  This is enough for
large directional comparisons, but not enough for strong per-case
significance claims.

\paragraph{Model variance.}
LLM outputs remain stochastic.  The value of the frozen protocol is that it
lets us observe large aggregate shifts despite that noise, but borderline
cases can still flip across rounds.  Future work should either increase $N$
or explicitly combine repeated runs with counterfactual comparisons.

\paragraph{Single model.}
The primary quantitative results use one model family
(\texttt{deepseek-v4-flash}).  Skill effects may differ across models, as
different models have different priors and instruction-following behavior.
Cross-model validation is left as future work.

\paragraph{Static analysis limitation.}
The firmware cases in the main studies are analyzed through binary-level
static reasoning without a live runtime environment.  Consequently, the
available evidence is weaker than sanitizer-backed confirmation in
open-source cases.  This likely lowers the precision of the feedback signal
and should be interpreted as a limitation of the present study, not as a
property of the framework in general.

\paragraph{Attribution precision.}
The feedback labels and target-aware miss tags are still heuristic
approximations, not causal estimators.  A run that lands on the correct
function under oracle guidance may still have been recoverable without that
skill, while a run marked as a miss may still have been partially helped by
skill-induced attention shifts.  Rigorous causal validation still requires
counterfactual ablation.

\paragraph{Gate evaluation.}
The 94 gate decisions were produced under multiple development-time gate
regimes.  Earlier decisions preserve only candidate-side validation
outcomes, whereas later decisions include baseline-aware metadata and
non-inferiority parameters.  This historical mixture is useful for
retrospective diagnosis but imperfect for a clean scientific comparison.  A
future study should evaluate gate variants from fixed repository snapshots
under a single benchmark protocol.

\paragraph{Answer leakage.}
Known-CVE benchmarks can leak through identifiers, artifact names,
comments, or pre-existing PoCs.  While our blind workspaces are generated
automatically and audited to exclude solution-bearing material, we cannot
guarantee zero leakage, especially for well-known CVEs.  The clean rerun
study itself is evidence that this threat is real rather than hypothetical.

\paragraph{Oracle validity.}
Case validators encode trusted expectations and may themselves be
incomplete.  A validator that passes does not guarantee the vulnerability
is fully understood; a validator that fails does not guarantee the analysis
is wrong.  We therefore report validator provenance and keep some cases as
supporting diagnostics rather than promoting them to the main evidence set.

\paragraph{Benchmark curation and invalidated cases.}
Part of the experimental contribution in this project is procedural rather
than purely numerical: several earlier case configurations had to be
reclassified, split out, or removed from the main denominator after we
found leakage, duplicate-family ambiguity, or misleading score inflation.
This means the benchmark is not a static artifact handed down in advance,
but an evolving research object that had to be cleaned before it could
support defensible claims.  We mitigate this risk by preserving supporting
records, explicitly separating primary and diagnostic evidence, and
reporting when a case remains useful for engineering diagnosis but not for
the main aggregate.

\section{Conclusion}
\label{sec:conclusion}

We studied the transition from blind vulnerability-analysis runs to live
reusable skills.  The main claim of the paper is not that autonomous skill
improvement is solved.  Rather, we show that this transition can be made
explicit enough to study: blind remote execution, local external
confirmation, run archival, candidate generation, and publication control
can be connected into a single operational loop.

The empirical picture is now clearer than in the earlier development stage.
After removing answer-bearing skill content, clean reruns show that oracle
performance remains substantially above weak/no controls rather than
collapsing to zero.  The larger frozen protocol strengthens that result:
oracle-skill reaches $46.7\%$ true hits versus $11.1\%$ and $6.7\%$, while
reducing decoy hits from 28/45 and 24/45 to 2/45.  At the same time, the
effect is not uniform.  Family-specific distinguishing skills help
F9K1122 dramatically but harm FH451, and several frozen cases still remain
unsolved even under oracle guidance.

We also show that publication control is the real bottleneck of the
current loop.  The system can already produce candidates, rerun them, and
publish some of them into the live library.  The harder problem is not
candidate generation but deciding whether a validated run contains knowledge
that truly deserves to persist.  Replay-only gates can enforce safety and
non-degradation, but they still do not by themselves prove causal skill
contribution.

This observation is relevant beyond vulnerability analysis.  If the
run-to-library transition is difficult even in a domain with executable
external confirmation, it is likely harder in agent settings where such
confirmation is weak or absent.  Our current framework suggests a pragmatic
stance: external evidence can discipline skill evolution, but it does not
remove the need for explicit counterfactual comparison, publication policy,
and retention policy.

Future work should proceed along four directions: (1) extend the frozen
protocol to more case families and more rounds; (2) add counterfactual
ablation or retention-stage gates that better separate improvement from
redundancy; (3) broaden the use of runtime-backed evidence beyond current
firmware cases; and (4) compare server-side inline injection with
catalog-style selection under the same controlled benchmark.
\section*{Acknowledgments}
Acknowledgments are omitted in the current working draft.

\bibliographystyle{elsarticle-num}
\bibliography{references}

\end{document}
